\documentclass[11pt,a4paper]{article}
\usepackage[T1]{fontenc}
\usepackage[utf8]{inputenc}
\usepackage[english]{babel}
\usepackage{amsmath,amssymb}
\usepackage[margin=2cm]{geometry}
\setlength{\parindent}{0pt}
\setlength{\emergencystretch}{3em}
\allowdisplaybreaks

\title{About \texttt{false\_index\_ranges.tex}}
\author{}
\date{}
\begin{document}
\maketitle
This note explains what the test data in \texttt{false\_index\_ranges.tex} is used for. That file is input for the engine in \texttt{main.c}; it is read by the test \texttt{test\_index\_identities.c} in this folder.

\begin{enumerate}
\item \textbf{Purpose.} The file holds 16 false identities with Greek indices or with values that start from 0. It is a regression test against false proofs: the engine must declare each one false, with a counterexample certified by the independent checker.

\item \textbf{How the statements were built.} Each one is a true identity of \texttt{index\_ranges.tex} with one small change: a trace counted over one value less, a numeric index moved by one, the time component read as the last one instead of the first, a sign of the metric, the value of $\varepsilon$ with the indices in the other position. Their falsity was confirmed by the oracle of \texttt{tools/fuzz/fuzz\_tensors.py} in its \texttt{data} mode.

\item \textbf{What the engine needs in order to pass.} The same reading as for the true identities, and a point where the two sides differ, written with the names and the values of the indices of the statement. The test also alters each certificate and requires the checker to reject it.

\item \textbf{Format expected by the test.} One \texttt{\textbackslash item} per statement, written as \texttt{\$\textbackslash displaystyle ...\$}. The values of the indices are declared by the quantifiers, $\forall \mu \in \{0, 1, 2, 3\}$, by a summation sign, $\sum_{\mu=0}^{3}$, by a hypothesis, $\mu, \nu \in \{0, 1, 2, 3\}$, or by the declared value of $\varepsilon$, $\varepsilon^{0123} = 1$.

\item \textbf{How to run it.} From the top folder of the repository:
\begin{verbatim}
cc -std=c11 -O2 -Wall -Wextra -Werror -o /tmp/test_index_identities \
    tests/tensor_calculus/test_index_identities.c && /tmp/test_index_identities
\end{verbatim}

\item \textbf{Expected result.} The lines \texttt{PASS: 16 of 16 statements refuted, 16 with a certificate the checker accepts} and \texttt{PASS: 32 of 32 altered certificates rejected for the right reason}.

\end{enumerate}
\end{document}
